\section{Behaviour as \texorpdfstring{$p\to\tfrac12^-$}{p to 1/2}}
\label{a:sec:asymptotic}

The series \cref{a:eq:Aseries} also delivers the near-critical asymptotic, and does
so more transparently than the matching argument one is first tempted to try.

Fix $\varepsilon$ small. The summand $r^n/(2-S_n)$ involves two decays on very
different scales. The geometric factor $r^n = (1-\varepsilon)^n$ decays on the
scale $n\sim1/\varepsilon$, which is large. The survival probability $S_n$, by
contrast, is close to its critical behaviour $S_n\sim2/n$ over that whole range,
and so has already fallen to near zero after $O(1)$ generations. For the
overwhelming majority of the terms that carry weight, therefore, $S_n\approx0$ and
the summand is approximately $r^n/2$. Summing,
\begin{equation*}
\sum_{n=0}^{\infty}\frac{r^n}{2-S_n}
\;\approx\;
\frac12\sum_{n=0}^{\infty}r^n = \frac{1}{2\varepsilon},
\end{equation*}
so $1/A(p)\sim1/(2\varepsilon)$ and
\begin{equation}
\label{a:eq:Aasym}
\boxed{\;A(p)\sim2\varepsilon = 2(1-2p)\qquad\text{as }p\to\tfrac12^-.\;}
\end{equation}
Equivalently $1/A(p)\sim1/\bigl(2(1-2p)\bigr)$: the quasi-stationary mean diverges
like the reciprocal of the distance from criticality. In particular $A(p)$ meets
the axis at $p=\tfrac12$ with gradient $-4$, a fact used repeatedly below as a test
of candidate formulae.

\begin{theorem}[Near-critical survival amplitude]
\label{a:thm:A-near-critical}
As $p\uparrow\tfrac12$, with $\varepsilon=1-2p\downarrow0$,
\begin{equation}
\label{a:eq:A-near-critical}
A(p)=2\varepsilon
\left[1+O\!\left(\varepsilon\log\frac1\varepsilon\right)\right].
\end{equation}
In particular $A(p)\sim2(1-2p)$ and $1/A(p)\sim1/(2(1-2p))$.
\end{theorem}

\noindent
The series argument above yields the leading asymptotic. A fully rigorous bound on
the remainder proceeds by decomposing $1/(2-S_n)=\tfrac12+S_n/(2(2-S_n))$ and
controlling the error series by the critical bound $S_n\le2/(n+2)$, which produces
an $O(\log(1/\varepsilon))$ remainder before inversion
\cite{Athreya1972}; the details match the argument recorded in the engineering
source for this chapter. The empirical fit
$1-A/(2\varepsilon)\approx\varepsilon(\log(1/\varepsilon)+0.78)$ is consistent with
the order but is not used subsequently.



The same result can be reached through the continuum equation of
\ChM, and it is worth seeing why that route is the worse one. The
exact difference is
\begin{equation*}
S_{n+1}-S_n = (2p-1)S_n - pS_n^2
= -\varepsilon S_n - \tfrac12S_n^2 + \tfrac{\varepsilon}{2}S_n^2 ,
\end{equation*}
and dropping the $O(\varepsilon S_n^2)$ term gives
$\dd S/\dd n = -\varepsilon S-\tfrac12S^2$. Setting $u=1/S$ linearises this to
$u' = \varepsilon u+\tfrac12$, with solution
\begin{equation*}
u(n) = C\ee^{\varepsilon n} - \frac{1}{2\varepsilon},
\qquad
S(n) = \frac{1}{C\ee^{\varepsilon n}-\frac{1}{2\varepsilon}} .
\end{equation*}
Matching $S_n\sim A\ee^{-\varepsilon n}$ for small $\varepsilon$ identifies
$A=1/C$, and matching to the critical decay $S_n\sim2/n$ in the appropriate double
limit fixes $C\sim1/(2\varepsilon)$, recovering \cref{a:eq:Aasym}. Both steps are
heuristic, and the second is delicate; the series argument is preferable because it
needs no second matching at all.

\begin{remark}[Rate of approach]
\label{a:rem:rate}
Numerically, $A(p)/2\varepsilon$ approaches $1$ rather slowly. Evaluating
\cref{a:eq:prodFormula} in high precision gives $A/2\varepsilon = 0.90987$,
$0.98612$, $0.99814$, $0.99977$ and $0.999972$ at $\varepsilon=2\times10^{-2}$ down
to $2\times10^{-6}$. Fitting the deficit gives
$1-A/2\varepsilon\approx\varepsilon\bigl(\ln(1/\varepsilon)+c\bigr)$ with
$c\approx0.78$, so
\begin{equation*}
A(p) = 2\varepsilon\Bigl[1+O\bigl(\varepsilon\ln(1/\varepsilon)\bigr)\Bigr] .
\end{equation*}
The logarithmic correction is why \cref{a:eq:Aasym} is a good practical surrogate
only very close to criticality, and it is what fixes the crossover point chosen in
\cref{a:sec:practical}.
\end{remark}

